% GENERATED FILE. Edit canonical lessons, then run npm run generate:books.
% canonical-source-hash: fnv1a64:55985f6d0328c084
% canonical-lessons: JA-C225-oboeru, JA-C225-renshuu, JA-C225-fukushuu, JA-C225-shiraberu, JA-C225-jisho

\chapter{Practice and Revision}
\label{ch:ja-practice-and-revision}
\hlchaptermodality{225}

\begin{chapteropening}
\textbf{By the end of this chapter:} \emph{I can say how I study and practise.}

Complete the last lesson of chapter 225: I can say how I study and practise.
\end{chapteropening}

\section[\texorpdfstring{\ja{おぼえる} (oboeru)}{おぼえる (oboeru)}]{\ja{おぼえる} (oboeru) — to memorize; to learn by heart}

\label{lesson:JA-C225-oboeru}



\begin{quote}
Before the new one: say the Japanese for to teach, then the Japanese for to take lessons.
\end{quote}

\subsection*{You'll want to know: \ja{おぼえる}}
\textbf{\ja{おぼえる}} — \emph{oboeru} — \textquotedblleft{}to memorize; to learn by heart\textquotedblright{}.

\begin{quote}
\textbf{\ja{せんせいは} \ja{がくせいの} \ja{なまえを} \ja{おぼえます。}} — \emph{sensei wa gakusei no namae o oboemasu} — \textquotedblleft{}The teacher learns the students' names by heart.\textquotedblright{}
\end{quote}

Polite \textbf{\ja{おぼえます}}, \emph{oboemasu}; an ichidan verb.

\begin{cousinweb}
\textbf{\ja{おぼえる}} (\emph{oboeru}) is a native Japanese word (wago), inherited rather than borrowed from Chinese; its kanji was attached later as a kun'yomi (JMdict, KANJIDIC2).
\end{cousinweb}

\begin{grammarlens}[title={What you've built}]
One more word for everyday life at A2.
\end{grammarlens}

\subsection*{Your turn}
\begin{itemize}
\raggedright
  \item \emph{Say it:} \emph{oboeru}
  \item \emph{Say it:} \emph{oboeru}, once more
  \item \emph{Say it:} \emph{sensei wa gakusei no namae o oboemasu} — The teacher learns the students' names by heart.
  \item \emph{Recall:} say the Japanese for to teach, then the Japanese for to take lessons, then say \emph{oboeru} again
\end{itemize}

\begin{tcolorbox}[breakable,colback=teal!4,colframe=teal!35!black,title={Before you move on}]
What does \ja{おぼえる} mean? (\textquotedblleft{}To memorize; to learn by heart\textquotedblright{}.) Say it once more.
\end{tcolorbox}
\section[\texorpdfstring{\ja{れんしゅう} (renshū)}{れんしゅう (renshū)}]{\ja{れんしゅう} (renshū) — practice; training}

\label{lesson:JA-C225-renshuu}



\begin{quote}
Before the new one: say the Japanese for to take lessons, then the Japanese for to memorize.
\end{quote}

\subsection*{You'll want to know: \ja{れんしゅう}}
\textbf{\ja{れんしゅう}} — \emph{renshū} — \textquotedblleft{}practice; training\textquotedblright{}.

\begin{quote}
\textbf{\ja{まいあさ} \ja{うたの} \ja{れんしゅうを} \ja{します。}} — \emph{maiasa uta no renshū o shimasu} — \textquotedblleft{}I do singing practice every morning.\textquotedblright{}
\end{quote}

\begin{cousinweb}
\textbf{\ja{れんしゅう}} (\emph{renshū}) is Sino-Japanese (kango): two Chinese characters read with their on'yomi, readings borrowed from Chinese. They are \emph{ren}, \textquotedblleft{}drill, train\textquotedblright{} (Middle Chinese \emph{lenH}) and \emph{shū}, \textquotedblleft{}learn\textquotedblright{} (\emph{zip}) (KANJIDIC2; Middle Chinese in Baxter's transcription).
\end{cousinweb}

\begin{grammarlens}[title={What you've built}]
One more word for everyday life at A2.
\end{grammarlens}

\subsection*{Your turn}
\begin{itemize}
\raggedright
  \item \emph{Say it:} \emph{renshū}
  \item \emph{Say it:} \emph{renshū}, once more
  \item \emph{Say it:} \emph{maiasa uta no renshū o shimasu} — I do singing practice every morning.
  \item \emph{Recall:} say the Japanese for to take lessons, then the Japanese for to memorize, then say \emph{renshū} again
\end{itemize}

\begin{tcolorbox}[breakable,colback=teal!4,colframe=teal!35!black,title={Before you move on}]
What does \ja{れんしゅう} mean? (\textquotedblleft{}Practice; training\textquotedblright{}.) Say it once more.
\end{tcolorbox}
\section[\texorpdfstring{\ja{ふくしゅう} (fukushū)}{ふくしゅう (fukushū)}]{\ja{ふくしゅう} (fukushū) — revision; going over lessons}

\label{lesson:JA-C225-fukushuu}



\begin{quote}
Before the new one: say the Japanese for to memorize, then the Japanese for practice.
\end{quote}

\subsection*{You'll want to know: \ja{ふくしゅう}}
\textbf{\ja{ふくしゅう}} — \emph{fukushū} — \textquotedblleft{}revision; going over lessons\textquotedblright{}.

\begin{quote}
\textbf{\ja{まいにち} \ja{うちで} \ja{ふくしゅうを} \ja{します。}} — \emph{mainichi uchi de fukushū o shimasu} — \textquotedblleft{}I do my revision at home every day.\textquotedblright{}
\end{quote}

\begin{cousinweb}
\textbf{\ja{ふくしゅう}} (\emph{fukushū}) is Sino-Japanese (kango): two Chinese characters read with their on'yomi, readings borrowed from Chinese. They are \emph{fuku}, \textquotedblleft{}go back, again\textquotedblright{} (Middle Chinese \emph{bjuwk}) and \emph{shū}, \textquotedblleft{}learn\textquotedblright{} (\emph{zip}) (KANJIDIC2; Middle Chinese in Baxter's transcription).
\end{cousinweb}

\begin{grammarlens}[title={What you've built}]
One more word for everyday life at A2.
\end{grammarlens}

\subsection*{Your turn}
\begin{itemize}
\raggedright
  \item \emph{Say it:} \emph{fukushū}
  \item \emph{Say it:} \emph{fukushū}, once more
  \item \emph{Say it:} \emph{mainichi uchi de fukushū o shimasu} — I do my revision at home every day.
  \item \emph{Recall:} say the Japanese for to memorize, then the Japanese for practice, then say \emph{fukushū} again
\end{itemize}

\begin{tcolorbox}[breakable,colback=teal!4,colframe=teal!35!black,title={Before you move on}]
What does \ja{ふくしゅう} mean? (\textquotedblleft{}Revision; going over lessons\textquotedblright{}.) Say it once more.
\end{tcolorbox}
\section[\texorpdfstring{\ja{しらべる} (shiraberu)}{しらべる (shiraberu)}]{\ja{しらべる} (shiraberu) — to look up; to check}

\label{lesson:JA-C225-shiraberu}



\begin{quote}
Before the new one: say the Japanese for practice, then the Japanese for revision.
\end{quote}

\subsection*{You'll want to know: \ja{しらべる}}
\textbf{\ja{しらべる}} — \emph{shiraberu} — \textquotedblleft{}to look up; to check\textquotedblright{}.

\begin{quote}
\textbf{\ja{としょかんで} \ja{こたえを} \ja{しらべます。}} — \emph{toshokan de kotae o shirabemasu} — \textquotedblleft{}I look up the answer at the library.\textquotedblright{}
\end{quote}

Polite \textbf{\ja{しらべます}}, \emph{shirabemasu}; an ichidan verb.

\begin{cousinweb}
\textbf{\ja{しらべる}} (\emph{shiraberu}) is a native Japanese word (wago), inherited rather than borrowed from Chinese; its kanji was attached later as a kun'yomi (JMdict, KANJIDIC2).
\end{cousinweb}

\begin{grammarlens}[title={What you've built}]
One more word for everyday life at A2.
\end{grammarlens}

\subsection*{Your turn}
\begin{itemize}
\raggedright
  \item \emph{Say it:} \emph{shiraberu}
  \item \emph{Say it:} \emph{shiraberu}, once more
  \item \emph{Say it:} \emph{toshokan de kotae o shirabemasu} — I look up the answer at the library.
  \item \emph{Recall:} say the Japanese for practice, then the Japanese for revision, then say \emph{shiraberu} again
\end{itemize}

\begin{tcolorbox}[breakable,colback=teal!4,colframe=teal!35!black,title={Before you move on}]
What does \ja{しらべる} mean? (\textquotedblleft{}To look up; to check\textquotedblright{}.) Say it once more.
\end{tcolorbox}
\section[\texorpdfstring{\ja{じしょ} (jisho)}{じしょ (jisho)}]{\ja{じしょ} (jisho) — a dictionary}

\label{lesson:JA-C225-jisho}



\begin{quote}
Before the new one: say the Japanese for revision, then the Japanese for to look up.
\end{quote}

\subsection*{You'll want to know: \ja{じしょ}}
\textbf{\ja{じしょ}} — \emph{jisho} — \textquotedblleft{}a dictionary\textquotedblright{}.

\begin{quote}
\textbf{\ja{がくせいは} \ja{いつも} \ja{じしょを} \ja{つかいます。}} — \emph{gakusei wa itsumo jisho o tsukaimasu} — \textquotedblleft{}The students always use a dictionary.\textquotedblright{}
\end{quote}

\begin{cousinweb}
\textbf{\ja{じしょ}} (\emph{jisho}) is Sino-Japanese (kango): two Chinese characters read with their on'yomi, readings borrowed from Chinese. They are \emph{ji}, \textquotedblleft{}words\textquotedblright{} (Middle Chinese \emph{zi}) and \emph{sho}, \textquotedblleft{}write\textquotedblright{} (\emph{syo}) (KANJIDIC2; Middle Chinese in Baxter's transcription).
\end{cousinweb}

\begin{grammarlens}[title={What you've built}]
One more word for everyday life at A2.
\end{grammarlens}

\subsection*{Your turn}
\begin{itemize}
\raggedright
  \item \emph{Say it:} \emph{jisho}
  \item \emph{Say it:} \emph{jisho}, once more
  \item \emph{Say it:} \emph{gakusei wa itsumo jisho o tsukaimasu} — The students always use a dictionary.
  \item \emph{Recall:} say the Japanese for revision, then the Japanese for to look up, then say \emph{jisho} again
\end{itemize}

\begin{tcolorbox}[breakable,colback=teal!4,colframe=teal!35!black,title={Before you move on}]
What does \ja{じしょ} mean? (\textquotedblleft{}A dictionary\textquotedblright{}.) Say it once more.
\end{tcolorbox}
